%!TEX root = ../sztuthesis_main.tex
% 设置中文摘要

\phantomsection
\addcontentsline{toc}{section}{摘要}

\keywordscn{在线考试；远程监考；C/S架构；Electron框架；流媒体；信息安全}
\categorycn{TP391}
\begin{abstractcn}
随着在线教育和远程考试的普及，保障考试公平性与安全性成为重要挑战。传统监考方式
难以适应线上场景，现有远程监控工具亦存在不足。针对这些问题，本设计旨在开发一个
基于C/S架构的安全考试信息收发与监控系统。客户端采用Electron框架构建，提供图形
化配置界面，负责实时捕获考生屏幕与摄像头画面、采集系统信息（如MAC地址、IP地址、
CPU使用率、运行进程等）、并通过FFmpeg和Node-Media-Server进行视频流处理与本
地分发。服务端基于Golang和Gin框架开发，负责管理客户端连接与身份验证、显示客户
端连接状态、分发视频流地址、
通过WebSocket实现与客户端的实时双向通信（如下发指令、接收状态信息、发送可设置
优先级的通知并记录历史）。系统还关注了安全性，采用加密存储本地配置、安全分发服务
器访问凭证、Token认证等机制，并设计了禁止随意退出、单实例运行等防作弊措施。
通过在Windows和macOS环境下进行功能、性能、
兼容性、安全性及弱网测试，结果表明该系统功能完备、性能稳定、兼容性良好、安全性高，
能够有效满足上机考试或竞赛的实时监控需求，提升监考效率与考试公平性。

% 图X幅，表X个，参考文献X篇（四号宋体）

\end{abstractcn}